% Seccion: Fundamentos de dinámica de fluidos computacional (CFD)
\section{Fundamentos de dinámica de fluidos computacional (CFD)}
\label{sec:cap02_fundamentos_cfd}

La Dinámica de Fluidos Computacional (CFD) discretiza las ecuaciones continuas en un conjunto de ecuaciones algebraicas.
El método de volúmenes finitos (FVM) es el estándar de oro en resolvedores de ingeniería térmica debido a su conservación conservativa estricta.
Las ecuaciones de Navier-Stokes promediadas por Reynolds (RANS) resuelven el flujo medio y modelan la turbulencia a través de aproximaciones.
Los modelos de turbulencia de dos ecuaciones, como el k-épsilon y el SST k-omega, son ampliamente utilizados por su estabilidad y precisión.
El modelo SST k-omega de Menter destaca al combinar la precisión en la capa límite cercana a la pared con la robustez del k-épsilon en el canal libre.
